% TODO make hw stylesheet

\documentclass[10pt]{article}
\usepackage[letterpaper,top=1in,left=1in,right=1in]{geometry}
\usepackage[dvipsnames,svgnames]{xcolor}
\usepackage[shortlabels]{enumitem}
\usepackage[framemethod=TikZ]{mdframed}
\usepackage{amsmath,amssymb,amsthm}
\usepackage[colorlinks]{hyperref}
\usepackage{multicol}
\usepackage{tabularx}
\usepackage{bbm}
\usepackage{tikz-cd}

\hypersetup{
    colorlinks=true,
    linkcolor=blue,
    filecolor=magenta,
    urlcolor=magenta,
    pdftitle={MAT 535 HW}
}

\usepackage{mathtools}
\usepackage{thmtools}
\usepackage[textsize=footnotesize]{todonotes}
\usepackage{titling}
\usepackage{cancel}

\reversemarginpar
\mdfdefinestyle{mdbluebox}{roundcorner=10pt,innerbottommargin=9pt,
linecolor=blue,backgroundcolor=TealBlue!5,}
\declaretheoremstyle[headfont=\sffamily\bfseries\color{MidnightBlue},
mdframed={style=mdbluebox},]{thmbluebox}

\mdfdefinestyle{mdbloobox}{frametitlefont=\bfseries,innerbottommargin=8pt,
nobreak=true,backgroundcolor=TealBlue!5,linecolor=blue,}
\declaretheoremstyle[headfont=\bfseries\color{MidnightBlue},
mdframed={style=mdbloobox},headpunct={\\[3pt]},postheadspace=0pt,]{thmbloobox}

\mdfdefinestyle{mdredbox}{frametitlefont=\bfseries,innerbottommargin=8pt,
nobreak=true,backgroundcolor=Salmon!5,linecolor=RawSienna,}
\declaretheoremstyle[headfont=\bfseries\color{RawSienna},
mdframed={style=mdredbox},headpunct={\\[3pt]},postheadspace=0pt,]{thmredbox}

\mdfdefinestyle{mdgreenbox}{linecolor=ForestGreen,backgroundcolor=ForestGreen!5,
linewidth=2pt,rightline=false,leftline=true,topline=false,bottomline=false,}
\declaretheoremstyle[headfont=\bfseries\sffamily\color{ForestGreen!70!black},
mdframed={style=mdgreenbox},headpunct={ --- },]{thmgreenbox}

\mdfdefinestyle{mdorangebox}{linecolor=RedOrange,backgroundcolor=RedOrange!5,
linewidth=2pt,rightline=false,leftline=true,topline=false,bottomline=false,}
\declaretheoremstyle[headfont=\bfseries\sffamily\color{RedOrange!70!black},
mdframed={style=mdorangebox},headpunct={ --- },]{thmorangebox}

\mdfdefinestyle{mdblackbox}{linecolor=black,backgroundcolor=RedViolet!5!gray!5,
linewidth=3pt,rightline=false,leftline=true,topline=false,bottomline=false,}
\declaretheoremstyle[mdframed={style=mdblackbox}]{thmblackbox}

\declaretheoremstyle[spaceabove=6pt,spacebelow=6pt]{boldhead}
\declaretheoremstyle[spaceabove=6pt,spacebelow=6pt,headfont=\normalfont\itshape]{italichead}

\declaretheorem[style=boldhead,name=Problem]{problem}
\declaretheorem[style=boldhead,name=Problem,numbered=no]{problem*}
\declaretheorem[style=boldhead,name=Setup]{setup} % for config geo
\declaretheorem[style=boldhead,name=Example,sibling=problem]{example}
\declaretheorem[style=boldhead,name=$\bigstar$ Problem,sibling=problem]{niceprob}
\declaretheorem[style=boldhead,name=Theorem,sibling=problem]{theorem}
\declaretheorem[style=boldhead,name=Corollary,sibling=theorem]{corollary}
\declaretheorem[style=boldhead,name=Proposition,sibling=theorem]{prop}
\declaretheorem[style=boldhead,name=Lemma,numberwithin=problem]{lemma}
\declaretheorem[style=boldhead,name=Theorem,numbered=no]{theorem*}
\declaretheorem[style=boldhead,name=Lemma,numbered=no]{lemma*}
\declaretheorem[style=boldhead,name=Claim,numberwithin=problem]{claim}
\declaretheorem[style=boldhead,name=Claim,numbered=no]{claim*}
\declaretheorem[style=boldhead,name=Remark,numberwithin=problem]{remark}
\declaretheorem[style=boldhead,name=Remark,numbered=no]{remark*}
\declaretheorem[style=boldhead,name=Definition,sibling=theorem]{definition}
\declaretheorem[style=boldhead,name=Definition,numbered=no]{definition*}

\providecommand{\ol}{\overline}
\providecommand{\tensor}{\otimes}
\renewcommand{\hom}{\operatorname{Hom}}
\providecommand{\tor}{\operatorname{Tor}}
\providecommand{\ceil}[1]{\left\lceil #1 \right\rceil}
\providecommand{\floor}[1]{\left\lfloor #1 \right\rfloor}
\newcommand{\dd}{\,\mathrm{d}}
\providecommand{\ul}{\underline}
\providecommand{\eps}{\varepsilon}
\providecommand{\half}{\frac{1}{2}}
\providecommand{\CC}{\mathbb C}
\providecommand{\FF}{\mathbb F}
\providecommand{\II}{\mathbb I}
\providecommand{\NN}{\mathbb N}
\providecommand{\QQ}{\mathbb Q}
\providecommand{\RR}{\mathbb R}
\providecommand{\ZZ}{\mathbb Z}
\providecommand{\ind}{\mathbbm 1}
\providecommand{\dg}{^\circ}
\providecommand{\ii}{\item}
\providecommand{\setdiff}{\smallsetminus}
\providecommand{\alert}[1]{{\sffamily\textbf{\textcolor{blue}{#1}}}}
\providecommand{\inv}{^{-1}}
\providecommand{\mb}[1]{\mathbf{#1}}
\newcommand{\what}{\widehat}

\newenvironment{soln}{\begin{proof}[Solution]}{\end{proof}}

\providecommand{\pow}{\operatorname{Pow}}
\providecommand{\vocab}[1]{{\textbf{\textcolor{ForestGreen}{#1}}}}
\providecommand{\lcm}{\operatorname{lcm}}
\providecommand{\abs}[1]{\left\lvert #1\right\rvert}
\providecommand{\smabs}[1]{\lvert #1\rvert}
\providecommand{\innerprod}[1]{\langle\!\langle #1\rangle\!\rangle}
\providecommand{\norm}[1]{\left\| #1\right\|}
\providecommand{\smnorm}[1]{\| #1\|}
\providecommand{\gr}{\operatorname{gr}}

\usepackage{mathrsfs}

% place title at top right
\setlength{\droptitle}{-8ex}
\pretitle{\begin{flushright}\Large\bfseries}
\posttitle{\par\end{flushright}}
\preauthor{\begin{flushright}}
\postauthor{\end{flushright}}
\predate{\begin{flushright}}
\postdate{\end{flushright}}

\title{MAT 535 HW03}
\author{\large Aditya Pahuja}
\date{\large Due 02/19}
\begin{document}
\maketitle

\begin{problem}
    [\S17.1\#10a,d]
    \begin{enumerate}[(a)]
	\ii Prove that an arbitrary direct sum $\bigoplus_{i\in I} P_i$
	of projective modules $P_i$ is projective
	and that an arbitrary direct product $\prod_{j\in J}Q_j$
	of injective modules $Q_j$ is injective.
	\setcounter{enumi}{3}
	\ii Prove that $\tor^R_n(A,\bigoplus_{j\in J}B_j)
	\cong\bigoplus_{j\in J}\tor^R_n(A,B_j)$
	for any collection of $R$-modules $\{B_j\}_{j\in J}$.
    \end{enumerate}
\end{problem}

\begin{soln}
    (a) Let $\{P_i\}_{i\in I}$ be a collection of projective modules,
    $u\colon M\to N$ is an epimorphism, and
    $f\colon\bigoplus_{i\in I}P_i\to N$.
    Define  the canonical inclusion $q_i\colon P_i\to\bigoplus_{i\in I} P_i$
    and let $f_i\colon P_i\to N$ be $f_i = f\circ q_i$.
    Since the $P_i$ are projective, there exist $\tilde f_i\colon P_i\to M$
    making the diagram
    \begin{center}
        \begin{tikzcd}
	    & P_i\arrow[ld,dashed,"\exists\tilde f_i"']
	    \arrow[d,"f_i"]\\
	    M\arrow[r,"u"] & N\arrow[r] & 0
        \end{tikzcd}
    \end{center}
    commute. By the universal property of the coproduct,
    the maps $\tilde f_i$ factor through the $q_i$,
    i.e. there is a unique morphism
    $\tilde f\colon\bigoplus_{i\in I}P_i\to M$ such that
    $\tilde f\circ q_i = \tilde f_i$.
    Moreover, the universal property of the coproduct also says that $f$
    is uniquely characterized by the equations $f_i = f\circ q_i$.
    Since $(u\circ\tilde f)\circ q_i = f_i$ for all $i$,
    this implies that $u\circ\tilde f = f$, so $\tilde f$
    is the desired lift of $f$.

    The analogous result for injective modules follows
    from passing to the opposite category,
    or equivalently making the same argument but with
    the universal property of the product instead
    (since in that setting there would be maps $g_j\colon N\to Q_j$
    which need to be extended to maps $\tilde g_j\colon M\to Q_j$
    in a manner compatible with a given monomorphism $u\colon N\to M$).
    \\

    (d) Fix a projective resolution
    $\dots\to P_{j,2}\to P_{j,1}\to P_{j,0}\to B_j\to 0$ for each $j\in J$.
    Since direct sums of projective modules are projective by (a)
    and direct sums of exact sequences are exact, we obtain
    a projective resolution
    \begin{equation*}
	\cdots\to \bigoplus_{j\in J}P_{j,2}\to
	\bigoplus_{j\in J}P_{j,1}\to
	\bigoplus_{j\in J}P_{j,0}\to
	\bigoplus_{j\in J}B_j\to 0,
    \end{equation*}
    which, upon tensoring with $A$ and truncating, yields the complex
    \begin{equation*}
	\cdots\to
	A\tensor_R\bigoplus_{j\in J}P_{j,2}\to
	A\tensor_R\bigoplus_{j\in J}P_{j,1}\to
	A\tensor_R\bigoplus_{j\in J}P_{j,0}\to
	A\tensor_R\bigoplus_{j\in J}B_j\to 0.
    \end{equation*}
    The homology of this complex is exactly equal to
    $\tor_n^R(A,\bigoplus_{j\in J}B_j)$.
    On the other hand, noting that the tensor product commutes
    with direct sums, this complex is the direct sum of the complexes
    \begin{equation*}
	\cdots\to
	A\tensor_RP_{j,2}\to
	A\tensor_RP_{j,1}\to
	A\tensor_RP_{j,0}\to
	A\tensor_RB_j\to 0,
    \end{equation*}
    whose $n$th homology module is $\tor_n^R(A,B_j)$.
    The fact that homology commutes with direct sums then implies that
    the $n$th homology of
    \begin{equation*}
	\cdots\to
	\bigoplus_{j\in J}A\tensor_RP_{j,2}\to
	\bigoplus_{j\in J}A\tensor_RP_{j,1}\to
	\bigoplus_{j\in J}A\tensor_RP_{j,0}\to
	\bigoplus_{j\in J}A\tensor_RB_j\to 0
    \end{equation*}
    is exactly $\bigoplus_{j\in J}\tor_n^R(A,B_j)$.
\end{soln}

\begin{problem}
    [\S17.1\#15]
    \begin{enumerate}[(a)]
        \ii If $I$ is an ideal in $R$ and $M$ is an $R$-module,
	prove that $\tor_1^R(M,R/I)$ is isomorphic
	to the kernel of the map $M\otimes_RI\to M$
	that maps $m\otimes i$ to $mi$ for $i\in I$ and $m\in M$.
	\ii Prove that the $R$-module $M$ is flat
	if and only if $\tor_1^R(M,R/I) = 0$
	for every finitely generated ideal $I$ of $R$.
    \end{enumerate}
\end{problem}

\begin{soln}
    (a) The short exact sequence $0\to I\to R\to R/I\to 0$
    is transformed by $\tor_\bullet^R(M,-)$ into the long exact sequence
    \begin{equation*}
        \dots\to\tor_1^R(M,R)\to\tor_1^R(M,R/I)\to
	M\tensor_RI\to M\tensor_RR\to M\tensor_R R/I\to 0.
    \end{equation*}
    There is an isomorphism $M\tensor_RR\cong M$ by sending
    $m\tensor r\mapsto mr$ (since every $m\tensor r$
    can be expresed uniquely as $n\tensor 1$ for some $n\in M$),
    so the morphism $M\tensor_RI\to M\tensor_RR\to M$
    given by composing this isomorphism with the
    morphism in the exact sequence induced by the inclusion
    $I\hookrightarrow M$ is the composed map
    $m\tensor i\mapsto m\tensor i\mapsto mi$.

    Since $R$ is a free $R$-module, it is flat,
    so $\tor_1^R(M,R) = 0$. By exactness of the sequence above
    at $\tor_1^R(M,R/I)$, this means that $\tor_1^R(M,R/I)$ maps injectively
    into $M\otimes_R I$. Then, by exactness at $M\tensor_R I$,
    the kernel of the $M\tensor_RI\to M\tensor_RR\cong M$ map
    is equal to the image of the preceding map in the sequence,
    which we just established is injective,
    so this image (hence the kernel of the map in question)
    is isomorphic to $\tor_1^R(M,R/I)$, as desired.
    \\

    (b) If $M$ is flat then $\tor_1^R(M,N) = 0$ for all $N$.

    Conversely, suppose we are given that
    $\tor_1^R(M,R/I) = 0$ for every finitely generated ideal $I\subseteq R$;
    in the long exact sequence for $\tor_\bullet^R(M,-)$, this means
    \begin{equation*}
        0\to M\tensor_R I\to M\tensor_RR\to M\tensor_RR/I\to 0
    \end{equation*}
    is an exact sequence. Flatness then follows from $\S10.5\#25$.
\end{soln}

\begin{problem}
    [\S17.1\#16]
    Suppose $R$ is a PID and $A$ and $B$ are $R$-modules.
    If $t(B)$ denotes the torsion submodule of $B$, show that
    $\tor_1^R(A,t(B))\cong\tor_1^R(A,B)$ and deduce that
    $\tor_1^R(A,B)\cong\tor_1^R(t(A),t(B))$.
\end{problem}

\begin{soln}
    Let $B' = t(B)$ and $B'' = B/t(B)$.
    For any ideal $I = (a)$ of $R$, the map
    \begin{equation*}
        B''\tensor_RI\to B''\tensor_RR\cong B'',\quad
	m\tensor ra\mapsto mra
    \end{equation*}
    is injective, since if $m_1r_1a = m_2r_2a$ then
    $m_1r_1 - m_2r_2$ is annihilated by $a$,
    meaning it is zero because $B/t(B)$ is torsionfree,
    so $m_1\tensor r_1a = m_1r_1\tensor a
    = m_2r_2\tensor a = m_2\tensor r_2a$.
    Since the kernel of this map is zero, the long exact sequence for
    $\tor_\bullet^R(B'',-)$ applied to the short exact sequence
    $0\to I\to R\to R/I\to 0$ yields $\tor_1^R(B'',R/I) = 0$,
    hence $B''$ is flat by part (b) of the previous problem.
    This in turn means that $\tor_n^R(N,B'') = 0$ for all $R$-modules $N$,
    so the long exact sequence for $\tor_\bullet^R(A,-)$
    applied to $0\to B'\to B\to B''\to 0$ looks like
    \begin{equation*}
	\cdots\to\underbrace{\tor_{n+1}^R(A,B'')}_{=0}\to
	\tor_n^R(A,B')\stackrel{\cong}\to\tor_n^R(A,B)
	\to\underbrace{\tor_n^R(A,B'')}_{=0}\to\cdots
    \end{equation*}
    meaning $\tor_n^R(A,B)\cong\tor_n^R(A,t(B))\cong\tor_n^R(t(A),t(B))$
    for all $n \ge 1$.
\end{soln}

\pagebreak
\begin{problem}
    [\S17.1\#23]
    Let $D\inv R$ be the localization of the commutative ring $R$
    with respect to the multiplicative subset $D$ of $R$.
    Prove that localization commutes with $\tor$, i.e.
    \[
	D\inv\tor_n^R(A,B)\cong\tor_n^{D\inv R}(D\inv A,D\inv B)
    \]
    for all $R$-modules $A$ and $B$ and all $n\ge 0$.
\end{problem}

\begin{soln}
    By exercise \S17.1\#22 with $S = D\inv R$
    and $f$ the obvious inclusion of $R$ into $S$, we have that
    $\tor_n^{D\inv R}(D\inv A,D\inv B) = \tor_n^R(A,D\inv B)$.
    For a given projective resolution $\dots\to P_1\to P_0\to A\to 0$
    of $A$, $\tor_\bullet^R(A,D\inv B)$ is the homology of
    \begin{equation*}
        \dots\to P_1\tensor_R D\inv B
	\to P_0\tensor_R D\inv B
	\to 0.
    \end{equation*}
    On the other hand, $-\tensor_RD\inv R$ is an exact functor
    because $D\inv R$ is flat, so it preserves
    kernels and cokernels, hence commutes with homology.
    Since $D\inv B\cong B\tensor D\inv R$,
    we can rewrite the chain complex above as
    \begin{equation*}
        \dots\to(P_1\tensor_R B)\tensor_RD\inv R
	\to(P_0\tensor_R B)\tensor_RD\inv R
	\to 0,
    \end{equation*}
    and its homology is
    $\tor_\bullet^R(A,B)\tensor_RD\inv R
    = D\inv\tor_\bullet^R(A,B)$.
\end{soln}

\begin{problem}
    [\S17.1\#26]
    \begin{enumerate}[(a)]
        \ii Prove that every finitely generated module
	over a Noetherian ring $R$ is finitely presented.
	\ii Prove that an $R$-module $M$ is finitely presented
	and projective if and only if $M$ is a direct summand
	of $R^n$ for some integer $n\ge 1$.
    \end{enumerate}
\end{problem}

\begin{soln}
    (a) Let $M$ be a finitely generated $R$-module,
    say with $t$ generators. Then there is a surjective map
    $q\colon R^t\to M$ sending each basis element to a generator of $M$.
    The kernel of this map must be finitely generated
    because submodules of finitely generated modules
    over Noetherian rings are finitely generated.
    Supposing that $\ker q$ has $s$ generators, there is an exact sequence
    $R^s\to R^t\overset q\to M\to 0$ where the leftmost map
    sends each basis element of $R^s$ to one of the generators of $\ker q$;
    the existence of this exact sequence is what it means
    to be finitely presented.
    \\

    (b) If $M$ is a direct summand of $R^n$ then it is of course projective.
    The projection $R^n\to M$ shows that $M$ is finitely generated,
    since it is generated by the images of the basis elements of $R^n$;
    (a) then shows that $M$ is finitely presented.

    Conversely, suppose $M$ is finitely presented and projective.
    Then there is a positive integer $t$
    and a surjection $R^t\overset q\to M$ (this is part
    of the exact sequence we obtain by $M$ being finitely presented).
    Since $M$ is projective, the identity map $M\to M$ lifts
    to a section $\sigma\colon M\to R^t$ such that
    $\sigma\circ q = \operatorname{id}$; therefore,
    the short exact sequence
    \begin{center}
        \begin{tikzcd}
	    0\arrow[r]&\ker q\arrow[r] &
	    R^t\arrow[r,"q"] & M
	    \arrow[bend left=33, dashed]{l}{\sigma}\arrow[r]&0
        \end{tikzcd}
    \end{center}
    splits, by the splitting lemma.
\end{soln}










\end{document}
